\section{The symplectic object of a finite field, and its twisted operator algebra}
\label{sec:foundations}

Throughout the labbook $p$ denotes a prime, $m \geq 1$ an integer, $q := p^m$,
and $\kappa := \mathbb{F}_q$ the finite field of order $q$, with prime field
$\mathbb{F}_p \subseteq \kappa$. \textbf{The case $p = 2$ is in scope
everywhere in this campaign}: no definition or theorem below is stated only
for odd $p$ and silently assumed to extend; where a statement holds only for
$p=2$ or only for odd $p$ this is marked in the statement itself, never left
to a reader's guess. This section builds the ground floor of the
construction — the symplectic vector space attached to $\kappa$, the two
choices (a polarizing cocycle and an additive character) the quantum system
is built from, the twisted operator algebra those choices produce, and its
associated Heisenberg group — and proves the three structural facts that hold
\emph{uniformly} in the characteristic and for \emph{every} admissible choice
of cocycle. What is and is not canonical about the choices themselves is the
subject of later sections, in particular \S\ref{sec:polarizing-cocycle}.

\subsection{The symplectic space of a finite field}

\begin{definition}[The symplectic object attached to a finite field]
\label{def:symplectic-object}
For a finite field $\kappa$, put $V(\kappa) := \kappa \oplus \kappa$, and
define
\[
\omega : V(\kappa) \times V(\kappa) \to \kappa, \qquad
\omega\big((a,b),(a',b')\big) := ab' - a'b .
\]
\end{definition}

\begin{scope}
This is a stipulation: a vector space and a bilinear pairing written down by
formula. On its own it asserts nothing about $\omega$ — not that it is
$\kappa$-bilinear, not that it is alternating, not that it is nondegenerate.
Those are properties, established as Theorem~\ref{thm:wh-form} below. No
restriction on the characteristic is built into the definition; $p=2$ is
included on the same footing as every other prime.
\end{scope}

\provenance{D1}{---}{---}{---}

\subsection{Polarizing cocycles: the datum $\beta$, introduced}
\label{ssec:beta-introduced}

The Weyl operators of \S\ref{ssec:weyl-algebra} below are built not from
$\omega$ directly but from a bilinear map whose antisymmetrization is
$\omega$. Choosing one such map is a second datum of the construction, on
exactly the same footing as the choice of additive character
$\psi$ of \S\ref{ssec:phase-datum}. This subsection only introduces the
datum and its notation; the fact that it forms a torsor, that it is
immaterial at odd characteristic, and that it splits the construction into
two genuinely different isomorphism types at $p=2$ is the content of
Theorems~\ref{thm:wh-beta-a}--\ref{thm:wh-beta-type} in
\S\ref{sec:polarizing-cocycle}, and is not anticipated here.

\begin{definition}[Admissible polarizing cocycles, and the reference cocycle]
\label{def:admissible-cocycles}
Fix a finite field $\kappa$ and its symplectic object $(V(\kappa),\omega)$ of
Definition~\ref{def:symplectic-object}. The set of \emph{admissible
polarizing cocycles} is
\[
\mathrm{Adm}(\omega) := \big\{\, \beta : V(\kappa)\times V(\kappa) \to \kappa
\ \big|\ \beta \text{ is $\kappa$-bilinear and } \beta - \beta^T = \omega
\,\big\},
\]
where $\beta^T(v,v') := \beta(v',v)$. A \emph{polarizing cocycle} is a choice
of element $\beta \in \mathrm{Adm}(\omega)$, and it is a \textbf{datum of the
construction}, on the same footing as the character $\psi$ of
Definition~\ref{def:phase-datum} — not a convention fixed once and for all.
The map
\[
\beta_0\big((a,b),(a',b')\big) := ab'
\]
lies in $\mathrm{Adm}(\omega)$ and is called the \emph{reference cocycle}: a
label for a point of the torsor of Definition~\ref{def:admissible-cocycles},
not a canonical point of it (Theorem~\ref{thm:wh-beta-a}).
\end{definition}

\begin{scope}
Nothing here asserts that $\mathrm{Adm}(\omega)$ is nonempty, that it has any
particular size, that $\beta_0$ is a typical or a distinguished element of
it, or that the choice of $\beta$ is or is not immaterial to anything built
from it. All of that is claimed, and proved, later:
Theorems~\ref{thm:wh-beta-a}--\ref{thm:wh-beta-type} of
\S\ref{sec:polarizing-cocycle} are where this datum's structure is actually
established. In particular the tempting
odd-characteristic habit of writing $\omega/2$ for ``the'' polarizing cocycle
is not part of this definition and is not available at $p=2$, where division
by $2$ is division by $0$; \S\ref{sec:polarizing-cocycle} makes this
precise.
\end{scope}

\provenance{D2}{---}{---}{---}

\subsection{The phase datum}
\label{ssec:phase-datum}

\begin{definition}[Phase datum, and the trace-normalized family]
\label{def:phase-datum}
A \emph{phase datum} for $\kappa$ is a nontrivial additive character
$\psi : (\kappa,+) \to \mathbb{C}^\times$, i.e.\ a group homomorphism with
$\psi \not\equiv 1$. Write $X(\kappa)$ for the set of such $\psi$. Inside
$X(\kappa)$, the \emph{trace-normalized family} is defined as follows: for
$\kappa = \mathbb{F}_{p^m}$ put
\[
\mathrm{Tr}_{\kappa/\mathbb{F}_p}(z) := z + z^p + z^{p^2} + \cdots +
z^{p^{m-1}} ,
\]
and for $\zeta$ a primitive $p$-th root of unity in $\mathbb{C}$, put
\[
\psi_\zeta(z) := \zeta^{\,\mathrm{Tr}_{\kappa/\mathbb{F}_p}(z)} .
\]
\end{definition}

\begin{scope}
Two choices are named here and neither is suppressed: the character $\psi$
itself (an arbitrary element of $X(\kappa)$), and, when working inside the
distinguished family, the root of unity $\zeta$. The absolute trace
$\mathrm{Tr}_{\kappa/\mathbb{F}_p}$ is canonical — it is the same
$\mathbb{F}_p$-linear map regardless of any further choice — but $\psi_\zeta$
is canonical only once $\zeta$ is fixed, and \S\ref{sec:functoriality} shows
in detail how far a compatible choice of $\psi$ across a tower of fields can
and cannot be made.
\end{scope}

\provenance{D3}{---}{---}{---}

\subsection{Weyl operators, the twisted algebra, and the Heisenberg group}
\label{ssec:weyl-algebra}

\begin{definition}[Weyl operators and the twisted observable algebra]
\label{def:weyl-algebra}
Fix a phase datum $\psi$ as in Definition~\ref{def:phase-datum} and a
polarizing cocycle $\beta \in \mathrm{Adm}(\omega)$ as in
Definition~\ref{def:admissible-cocycles}. The \emph{twisted operator algebra}
$A_{\psi,\beta}(V)$ is the $\mathbb{C}$-vector space
\[
A_{\psi,\beta}(V) := \bigoplus_{v \in V(\kappa)} \mathbb{C}\cdot W_\beta(v)
\]
with the product fixed, once and for all, on basis elements by
\[
W_\beta(v)\, W_\beta(v') := \psi\big(\beta(v,v')\big)\, W_\beta(v+v') .
\]
\end{definition}

\begin{scope}
The ordering of the product on the right ($\psi(\beta(v,v'))$, not
$\psi(\beta(v',v))$) is fixed as part of the definition and is not
symmetrized; which of the two orderings a later stipulated model (such as
$M_0$ of \S\ref{sec:models}) actually realizes is a computation, not a
convention, and is checked explicitly there. $A_{\psi,\beta}(V)$ depends on
$\kappa$, $\psi$ and $\beta$ and on nothing else; that this dependence is
associative, unital and of dimension $q^2$ is Theorem~\ref{thm:wh-comm}(a)
below, not part of the stipulation.
\end{scope}

\provenance{D4}{---}{---}{---}

\begin{definition}[The Heisenberg group of a polarizing cocycle]
\label{def:heisenberg-group}
With $\beta \in \mathrm{Adm}(\omega)$ as above, the \emph{Heisenberg group of
$\beta$} is the set $H_\beta(\kappa) := \kappa \times V(\kappa)$ with product
\[
(t,v)(t',v') := \big(t+t'+\beta(v,v'),\ v+v'\big).
\]
\end{definition}

\begin{scope}
This construction uses no character $\psi$; it is a group attached to
$\kappa$ and $\beta$ alone. It does depend on $\beta$, and — this is the
content of \S\ref{sec:polarizing-cocycle} — at $p=2$ its isomorphism class as
an abstract group depends on more than the isomorphism class of the pair
$(V(\kappa),\omega)$: it depends on which point of the torsor
$\mathrm{Adm}(\omega)$ was chosen. Nothing about that dependence is claimed
or used here.
\end{scope}

\provenance{D5}{---}{---}{---}

\subsection{The form is bilinear, alternating, nondegenerate, and its
isometries}

\statusProved
\begin{theorem}[The symplectic form of a finite field, at every characteristic]
\label{thm:wh-form}
Let $\kappa$ be any finite field, $p=2$ included, with $V(\kappa)$ and
$\omega$ as in Definition~\ref{def:symplectic-object}. Then:
\begin{enumerate}[label=(\alph*)]
\item $\omega$ is $\kappa$-bilinear;
\item $\omega(v,v) = 0$ for every $v \in V(\kappa)$ — \emph{alternating} in
the strong sense, not merely antisymmetric;
\item $\omega$ is nondegenerate: if $\omega(v,\cdot) \equiv 0$ then $v=0$;
\item every $\beta \in \mathrm{Adm}(\omega)$ (Definition~\ref{def:admissible-cocycles})
is a $2$-cocycle for the abelian group $(V(\kappa),+)$, i.e.
$\beta(v,v')+\beta(v+v',v'') = \beta(v,v''+v)+\beta(v',v'')$ for all
$v,v',v''$ — equivalently $\beta(v,v')+\beta(v+v',v'')=\beta(v',v'')+\beta(v,v'+v'')$
— and $\beta - \beta^T = \omega$;
\item the group of $\kappa$-linear automorphisms $g$ of $V(\kappa)$ with
$\omega(gv,gv')=\omega(v,v')$ for all $v,v'$ is exactly $SL_2(\kappa)$;
\item the group of merely $\mathbb{F}_p$-linear automorphisms $g$ of
$V(\kappa)$ (i.e.\ additive maps that need not commute with multiplication by
$\kappa$) with $\omega(gv,gv')=\omega(v,v')$ for all $v,v'$ is \emph{also}
exactly $SL_2(\kappa)$ — the larger class of symmetries one might a priori
expect contains nothing beyond the $\kappa$-linear ones.
\end{enumerate}
\end{theorem}

\begin{scope}
Part (f) is a statement about the $\kappa$-valued form $\omega$ itself; it is
emphatically \emph{not} a statement about the $\mathbb{F}_p$-valued form
$\psi\circ\omega$ obtained by composing with a character, whose isometries
are the strictly larger group $Sp_{2m}(\mathbb{F}_p)$ once $m>1$ — that is
Theorem~\ref{thm:wh-symm} of \S\ref{sec:fp-symmetry}. Nothing here addresses
which $\beta\in\mathrm{Adm}(\omega)$ is chosen, or classifies
$\mathrm{Adm}(\omega)$ itself; that is \S\ref{sec:polarizing-cocycle}.
\end{scope}

\provenance{D1, D2, WH-FORM}{theory/wh-kappa.md \S1}{theory/checks/wh\_kappa\_check.py (C1,C2); theory/checks/fp\_symmetry.py (S1)}{theory/verdicts/wh-kappa-r1.md}

\begin{proof}
(a) Fix $v=(a,b)$, $v'=(a',b')$. Each argument of $\omega$ occurs exactly
once in the formula $ab'-a'b$, and $\kappa$ is commutative, so scaling either
argument by $c\in\kappa$ scales the value by $c$, and additivity in each slot
is distributivity of multiplication over addition in $\kappa$. This holds
verbatim at $p=2$: no step divides by $2$ or by anything else.

(b) $\omega\big((a,b),(a,b)\big) = ab-ab=0$, an identity in the commutative
ring $\kappa$ that cancels two literally equal terms, not two terms related
by a factor of $2$. It therefore holds at $p=2$ exactly as it holds
elsewhere. This is worth pausing on: antisymmetry alone, $\omega(v,v')=-\omega(v',v)$,
gives only $2\omega(v,v)=0$ on setting $v'=v$, which is vacuous once $2=0$.
So at $p=2$, (b) is a strictly stronger statement than antisymmetry and must
be — and here is — verified by the direct computation above, not inferred
from it.

(c) If $\omega(v,\cdot)\equiv 0$ for $v=(a,b)$ then in particular
$\omega\big((a,b),(0,1)\big)=a=0$ and $\omega\big((a,b),(1,0)\big)=-b=0$, so
$v=0$.

(d) $\beta-\beta^T=\omega$ is the definition of $\mathrm{Adm}(\omega)$
(Definition~\ref{def:admissible-cocycles}). For the cocycle identity: $\beta$
is bi-additive (it is $\kappa$-bilinear, hence in particular additive in each
argument), and for any bi-additive $\beta$ both sides of the cocycle identity
expand, term by term, to $\beta(v,v')+\beta(v,v'')+\beta(v',v'')$; the
identity holds without using anything about $\omega$ beyond bi-additivity of
$\beta$.

(e) Write $v=(a,b)$, $v'=(a',b')$ as the columns of a $2\times 2$ matrix over
$\kappa$; then $\omega(v,v')=ab'-a'b=\det[v\,|\,v']$ by direct expansion. For
$g\in GL_2(\kappa)$, $[gv\,|\,gv']=g[v\,|\,v']$ as matrices, and
multiplicativity of the determinant for $2\times 2$ matrices over a
commutative ring (itself a polynomial identity verified by direct expansion,
holding in every characteristic) gives $\omega(gv,gv')=\det(g)\cdot\omega(v,v')$.
Hence $\omega\circ g=\omega$ iff $(\det g - 1)\,\omega \equiv 0$; since
$\omega\big((1,0),(0,1)\big)=1\neq 0$, this holds iff $\det g = 1$, i.e.\ iff
$g\in SL_2(\kappa)$.

(f) Let $g:V(\kappa)\to V(\kappa)$ be additive (not assumed $\kappa$-linear)
with $\omega(gv,gu)=\omega(v,u)$ for all $v,u$.
\begin{itemize}
\item \emph{$g$ is bijective.} If $gv=0$ then $\omega(v,u)=\omega(gv,gu)=0$
for every $u$, so $v=0$ by (c); thus $g$ is injective, hence bijective since
$V(\kappa)$ is finite.
\item \emph{$g$ sends the standard basis to a $\kappa$-basis.} Let
$v_1=(1,0)$, $v_2=(0,1)$, so $\omega(v_1,v_2)=1$. Then
$\omega(gv_1,gv_2)=\omega(v_1,v_2)=1\neq 0$; were $gv_1,gv_2$
$\kappa$-dependent, bilinearity and (b) would force $\omega(gv_1,gv_2)=0$
(if, say, $gv_2=c\cdot gv_1$ then $\omega(gv_1,gv_2)=c\,\omega(gv_1,gv_1)=0$
by (a) and (b)). So $\{gv_1,gv_2\}$ is $\kappa$-independent, hence a basis of
the $2$-dimensional space $V(\kappa)$.
\item \emph{$g$ is $\kappa$-linear.} Fix $c\in\kappa$, $u\in V(\kappa)$, and
put $w := g(cu) - c\cdot g(u)$ (subtraction makes sense because $g$ is
additive, hence $\mathbb{Z}$-linear, even before $\kappa$-linearity is
known). For $i=1,2$,
\[
\omega(gv_i,w) = \omega\big(gv_i,g(cu)\big) - c\,\omega(gv_i,gu)
= \omega(v_i,cu) - c\,\omega(v_i,u) = c\,\omega(v_i,u)-c\,\omega(v_i,u) = 0,
\]
using the hypothesis $\omega(g\cdot,g\cdot)=\omega(\cdot,\cdot)$ for the
middle equality and $\kappa$-bilinearity (a) for the outer ones. Since
$\{gv_1,gv_2\}$ is a basis, $\omega(gv_1,\cdot)$ and $\omega(gv_2,\cdot)$ have
trivial common kernel only at $0$ by nondegeneracy (c) applied after a change
of basis, so $w=0$. Hence $g(cu)=c\cdot g(u)$ for all $c,u$: $g$ is
$\kappa$-linear.
\end{itemize}
By $\kappa$-linearity and (e), $\det g = 1$, so $g\in SL_2(\kappa)$; the
converse containment is immediate from (e). This closes (f).
\end{proof}

\begin{remark}
Part (f) forecloses a tempting shortcut: one might hope to characterize the
$\mathbb{F}_p$-linear isometries of $\omega$ as the ``$\omega$-self-adjoint
$\mathbb{F}_p$-endomorphisms,'' by analogy with the real symplectic case.
This is false at $p=2$: exhaustive enumeration finds $8$ and $64$
$\omega$-self-adjoint $\mathbb{F}_p$-endomorphisms of $V(\kappa)$ at $q=2,4$
respectively, against only $2$ and $4$ $\kappa$-scalars — the self-adjointness
route simply does not isolate $SL_2(\kappa)$ once $p=2$, and the direct
argument above is what closes (f) instead. Independently, complete
backtracking search over the $\mathbb{F}_p$-linear isometries of the
$\kappa$-valued form finds exactly $|SL_2(\kappa)| = 6,24,60,120,504,720$ of
them at $q=2,3,4,5,8,9$, matching an independent enumeration of $SL_2(\kappa)$
set-for-set at every one of these values (\texttt{theory\slash{}checks\slash{}fp\_symmetry.py},
gates S1 and S4).
\end{remark}

\subsection{Weyl relations, uniformly in the characteristic and for every
choice of cocycle}

\statusProved
\begin{theorem}[The Weyl commutation relations, for every admissible cocycle]
\label{thm:wh-comm}
Let $\kappa$ be any finite field, $\psi$ a phase datum
(Definition~\ref{def:phase-datum}), and $\beta \in \mathrm{Adm}(\omega)$ an
\emph{arbitrary} admissible polarizing cocycle
(Definition~\ref{def:admissible-cocycles}). Then, inside $A_{\psi,\beta}(V)$
of Definition~\ref{def:weyl-algebra}:
\begin{enumerate}[label=(\alph*)]
\item $A_{\psi,\beta}(V)$ is a unital associative $\mathbb{C}$-algebra of
dimension $q^2$, with unit $W_\beta(0)$;
\item every $W_\beta(v)$ is invertible, with
$W_\beta(v)^{-1} = \psi\big(\beta(v,v)\big)\, W_\beta(-v)$;
\item $W_\beta(v)\,W_\beta(v') = \psi\big(\omega(v,v')\big)\, W_\beta(v')\,W_\beta(v)$
for all $v,v'\in V(\kappa)$;
\item $W_\beta(u)\,W_\beta(v)\,W_\beta(u)^{-1} = \psi\big(\omega(u,v)\big)\,W_\beta(v)$
for all $u,v\in V(\kappa)$.
\end{enumerate}
\end{theorem}

\begin{scope}
No part of this theorem uses which admissible cocycle $\beta\in\mathrm{Adm}(\omega)$
was chosen, nor any property of $\beta$ beyond bi-additivity and the defining
relation $\beta-\beta^T=\omega$; in particular nothing here distinguishes the
reference cocycle $\beta_0$ from any other point of the torsor. No hypothesis
on the characteristic is used beyond what already went into
Theorem~\ref{thm:wh-form}. What \emph{is} $\beta$-dependent — the
isomorphism class of the group $H_\beta(\kappa)$ itself, at $p=2$ — is not
addressed here.
\end{scope}

\provenance{D1, D2, D3, D4, WH-FORM, WH-COMM}{theory/wh-kappa.md \S2}{theory/checks/wh\_kappa\_check.py (C3,C4)}{theory/verdicts/wh-kappa-r1.md}

\begin{proof}
(a) Associativity on basis elements amounts to
\[
\big(W_\beta(v)W_\beta(v')\big)W_\beta(v'')=W_\beta(v)\big(W_\beta(v')W_\beta(v'')\big),
\]
which unwinds to the $2$-cocycle identity for $\psi\circ\beta$; this holds
because $\beta$ is a $2$-cocycle (Theorem~\ref{thm:wh-form}(d)) and $\psi$ is
multiplicative. Since $\beta(0,v)=\beta(v,0)=0$ ($\kappa$-bilinearity), and
$\psi(0)=1$, $W_\beta(0)$ is a two-sided unit. The dimension is
$|V(\kappa)|=q^2$ by construction.

(b) By bi-additivity, $\beta(v,-v)=-\beta(v,v)$, so
$W_\beta(v)W_\beta(-v)=\psi\big(\beta(v,-v)\big)W_\beta(0)=\psi\big(-\beta(v,v)\big)\cdot 1$,
and symmetrically on the other side; since $\psi$ takes values in
$\mathbb{C}^\times$, this exhibits a two-sided inverse.

(c) By the defining product, $W_\beta(v)W_\beta(v')=\psi\big(\beta(v,v')\big)W_\beta(v+v')$
and $W_\beta(v')W_\beta(v)=\psi\big(\beta(v',v)\big)W_\beta(v+v')$. By (b),
$W_\beta(v+v')$ is invertible, so
\[
W_\beta(v)W_\beta(v') = \psi\big(\beta(v,v')-\beta(v',v)\big)\,W_\beta(v')W_\beta(v)
= \psi\big(\omega(v,v')\big)\,W_\beta(v')W_\beta(v),
\]
the last step by the defining relation $\beta-\beta^T=\omega$ of
Definition~\ref{def:admissible-cocycles}.

(d) Multiply (c) on the right by $W_\beta(v)^{-1}$, whose existence is (b).
\end{proof}

\begin{remark}[The centre of the Heisenberg group, and the algebra as a
group-algebra quotient]
The group $H_\beta(\kappa)$ of Definition~\ref{def:heisenberg-group} is
indeed a group: associativity is the cocycle identity of
Theorem~\ref{thm:wh-form}(d), the identity is $(0,0)$, and $(t,v)$ has
inverse $(-t+\beta(v,v),-v)$. Its centre is exactly $\kappa\times 0$: $(t,v)$
is central iff $\beta(v,v')-\beta(v',v)=0$ for every $v'$, i.e.\ iff
$\omega(v,\cdot)\equiv 0$, i.e.\ iff $v=0$ by
Theorem~\ref{thm:wh-form}(c). Consequently $A_{\psi,\beta}(V)$ is the
quotient of the group algebra $\mathbb{C}[H_\beta(\kappa)]$ by the two-sided
ideal identifying the central element $(t,0)$ with the scalar $\psi(t)$: the
algebra is a family over $\psi\in X(\kappa)$ at \emph{fixed} $\beta$, and the
group $H_\beta(\kappa)$ itself depends on $\beta$ — a dependence that,
Theorem~\ref{thm:wh-beta-type} shows, is not always trivial once $p=2$.
\end{remark}

\subsection{Simplicity of the twisted algebra}

\statusProved
\begin{theorem}[Central simplicity of the twisted algebra, with no polarization]
\label{thm:wh-alg}
With $\kappa$, $\psi$ and $\beta \in \mathrm{Adm}(\omega)$ as in
Theorem~\ref{thm:wh-comm}, the algebra $A_{\psi,\beta}(V)$ of
Definition~\ref{def:weyl-algebra} is simple, has centre $\mathbb{C}\cdot 1$,
and $\dim_{\mathbb{C}} A_{\psi,\beta}(V) = q^2$.
\end{theorem}

\begin{scope}
This theorem uses no choice of Lagrangian line, no module, and no property of
$\beta$ beyond Theorem~\ref{thm:wh-comm}; in particular it holds for every
$\beta\in\mathrm{Adm}(\omega)$, uniformly. It does not by itself identify
$A_{\psi,\beta}(V)$ with a matrix algebra $M_q(\mathbb{C})$ — that
identification exists (Theorem~\ref{thm:wh-alg-mat} of
\S\ref{sec:models}) but needs a polarization and is not canonical, whereas
simplicity and the dimension count here need none.
\end{scope}

\provenance{D4, WH-COMM, WH-ALG}{theory/wh-kappa.md \S3}{theory/checks/wh\_kappa\_check.py (C5,C6)}{theory/verdicts/wh-kappa-r1.md}

\begin{proof}
Define the \emph{Weyl average} $E : A_{\psi,\beta}(V) \to A_{\psi,\beta}(V)$
by
\[
E(x) := q^{-2}\sum_{u\in V(\kappa)} W_\beta(u)\, x\, W_\beta(u)^{-1}.
\]

\emph{Step 1: $E\big(\sum_v a_v W_\beta(v)\big) = a_0\cdot 1$.} By
Theorem~\ref{thm:wh-comm}(d), $W_\beta(u)W_\beta(v)W_\beta(u)^{-1} =
\psi(\omega(u,v))W_\beta(v)$, so
\[
E\Big(\sum_v a_v W_\beta(v)\Big) = \sum_v a_v \Big(q^{-2}\sum_u \psi\big(\omega(u,v)\big)\Big) W_\beta(v).
\]
Fix $v\neq 0$. By Theorem~\ref{thm:wh-form}(a),(c), $\omega(\cdot,v)$ is a
nonzero $\kappa$-linear map $V(\kappa)\to\kappa$, hence surjective (a nonzero
linear map between finite-dimensional vector spaces over the same field onto
a $1$-dimensional target is onto). So $u\mapsto\psi(\omega(u,v))$ is a
character of the additive group $V(\kappa)$ that is nontrivial (since $\psi$
is nontrivial and $\omega(\cdot,v)$ is onto $\kappa$). A nontrivial character
$\chi$ of a finite abelian group $G$ sums to zero: fixing $u_0$ with
$\chi(u_0)\neq 1$, the substitution $u\mapsto u+u_0$ permutes $G$, so
$S:=\sum_u\chi(u)$ satisfies $S=\chi(u_0)S$, forcing $S=0$ since
$\chi(u_0)\neq 1$. Applying this to $\chi = \psi(\omega(\cdot,v))$ gives
$\sum_u\psi(\omega(u,v))=0$ for $v\neq 0$; for $v=0$ the sum is
$q^2\cdot\psi(0)=q^2$. Hence $E(\sum_v a_vW_\beta(v)) = a_0\cdot W_\beta(0) =
a_0\cdot 1$, since $W_\beta(0)$ is the unit (Theorem~\ref{thm:wh-comm}(a)).

\emph{Step 2: $Z(A_{\psi,\beta}(V)) = \mathbb{C}\cdot 1$.} If $z$ is central
then $W_\beta(u)zW_\beta(u)^{-1}=z$ for every $u$, so $z=E(z)=a_0\cdot 1$ by
Step 1; conversely every scalar is central.

\emph{Step 3: simplicity.} Let $I$ be a nonzero two-sided ideal and pick
$0\neq x=\sum_v a_vW_\beta(v) \in I$ with some coefficient $a_{v_0}\neq 0$.
Then $y := W_\beta(v_0)^{-1}x \in I$ (as $I$ is a left ideal), and by
Theorem~\ref{thm:wh-comm}(b),(c) each $W_\beta(v_0)^{-1}W_\beta(v)$ is a
nonzero scalar multiple of $W_\beta(v-v_0)$; in particular the coefficient of
$W_\beta(0)$ in $y$ is a nonzero scalar multiple of $a_{v_0}$, hence nonzero.
Now $E(y)\in I$: $E$ is a $\mathbb{C}$-linear combination of the maps
$x\mapsto W_\beta(u)xW_\beta(u)^{-1}$, and $I$ is a two-sided ideal, so each
term $W_\beta(u)yW_\beta(u)^{-1}$ lies in $I$ and so does their average. By
Step 1, $E(y)=c\cdot 1$ with $c\neq 0$ (the coefficient of $W_\beta(0)$ in
$y$), so $1\in I$ and $I=A_{\psi,\beta}(V)$.

\emph{Step 4: dimension.} $\dim_{\mathbb{C}}A_{\psi,\beta}(V)=|V(\kappa)|=q^2$
directly from Definition~\ref{def:weyl-algebra}.
\end{proof}

\begin{remark}[Why the quantifier over $\beta$ can be widened after the
fact]
Theorems~\ref{thm:wh-comm} and \ref{thm:wh-alg} are stated for \emph{every}
$\beta\in\mathrm{Adm}(\omega)$, not merely the reference cocycle $\beta_0$.
This widening is legitimate precisely because no step of either proof used
more than the defining relation $\beta-\beta^T=\omega$ and bi-additivity: the
commutator computation in Theorem~\ref{thm:wh-comm}(c)--(d) depends on
$\beta$ only through $\omega=\beta-\beta^T$, and simplicity in
Theorem~\ref{thm:wh-alg} depends only on $\omega$ being nondegenerate. That
the choice of $\beta$ is nonetheless \emph{not} immaterial to everything
built from it — it is invisible to the pair $(A_{\psi,\beta}(V),\text{Weyl
frame})$ but not to the Heisenberg group $H_\beta(\kappa)$ itself once
$p=2$ — is the subject of \S\ref{sec:polarizing-cocycle}.
\end{remark}
